\documentclass[11pt,a4paper]{article}

\usepackage[margin=2.5cm]{geometry}
\usepackage{booktabs}
\usepackage{amsmath}
\usepackage{graphicx}
\usepackage{hyperref}
\usepackage{xcolor}
\usepackage{microtype}
\usepackage{parskip}
\usepackage{titlesec}
\usepackage{caption}
\usepackage{listings}
\usepackage{multirow}

\hypersetup{
    colorlinks=true,
    linkcolor=blue!60!black,
    citecolor=blue!60!black,
    urlcolor=blue!60!black
}

\titleformat{\section}{\large\bfseries}{\thesection}{1em}{}
\titleformat{\subsection}{\normalsize\bfseries}{\thesubsection}{1em}{}

\title{\textbf{Zero-Shot Event Extraction with Large Language Models:\\
A Baseline Study on MAVEN and WikiEvents}}

\author{
    Ross Feng\\
    \small University of Sheffield\\
    \small \texttt{acp25ck@sheffield.ac.uk}
}

\date{\today}

\begin{document}

\maketitle

\begin{abstract}
We evaluate whether large language models (LLMs) can perform event extraction with no task-specific training, using structured prompt engineering alone. Using \textbf{Qwen2.5-7B-Instruct} and \textbf{Llama-3.1-8B-Instruct}, we run zero-shot and few-shot inference on two benchmarks: MAVEN, a large-scale event detection dataset, and WikiEvents (NAACL 2021), a richer document-level event extraction benchmark. We compare six experimental conditions across three prompting strategies and two model families, evaluating with exact match, fuzzy string matching (Levenshtein distance), hierarchical partial credit, and micro/macro F1. Results show that constrained-label prompting and few-shot examples both improve type classification substantially, and that Llama-3.1-8B outperforms Qwen2.5-7B on all metrics despite identical parameter count. Multi-event sentence ambiguity and ontology misalignment remain the dominant failure modes across all conditions. All code, results, and slot-level evaluation templates are available at \url{https://github.com/rosscyking1115/event-extraction-llm-baseline}.
\end{abstract}

% ─────────────────────────────────────────────────────────────────────────────
\section{Introduction}

Event extraction is the task of identifying event triggers and their types in natural language text. Traditional approaches rely on supervised models trained on annotated corpora~\cite{ahn2006stages,li2013joint}. With the advent of capable instruction-tuned LLMs, an important question arises: \textit{can a general-purpose LLM perform event extraction without any task-specific fine-tuning?}

This report presents a systematic LLM baseline study for event extraction. We investigate three core questions:
\begin{enumerate}
    \item How accurately can instruction-tuned LLMs identify event triggers and types with no in-context examples?
    \item Does providing the model with a constrained label list or few-shot examples improve performance, and by how much?
    \item Does model family matter — do Qwen2.5-7B and Llama-3.1-8B exhibit different strengths and failure modes?
\end{enumerate}

We experiment on two datasets of increasing complexity: MAVEN~\cite{wang2020maven}, which focuses on event detection, and WikiEvents~\cite{li2021document}, which includes richer event argument structure and a three-level type hierarchy. Across six experimental conditions — including a rule-based majority-class lower bound — we find that constrained prompting and few-shot examples both provide meaningful improvements, that Llama outperforms Qwen despite identical scale, and that multi-event sentence ambiguity is the single largest failure mode regardless of model or prompting strategy.

% ─────────────────────────────────────────────────────────────────────────────
\section{Background and Related Work}

\subsection{Event Extraction}

Event extraction encompasses two subtasks: \textit{trigger detection} (identifying the word or phrase that most clearly expresses an event) and \textit{event classification} (assigning the trigger to a predefined event type from an ontology). Full event extraction additionally requires identifying \textit{event arguments} — the entities involved in each event and their roles.

Early systems used feature-based classifiers~\cite{ahn2006stages}. Neural approaches improved substantially~\cite{nguyen2015event,yang2019exploring}, and recent work uses generative models for end-to-end extraction~\cite{li2021document,lu2021text2event}.

\subsection{Zero-Shot LLMs for IE}

Several recent works have explored zero-shot and few-shot prompting for information extraction tasks~\cite{wei2023zero,ma2023large}. These studies generally find that LLMs perform well on high-level classification but struggle with fine-grained ontology alignment and exact span matching — findings that are consistent with our results.

% ─────────────────────────────────────────────────────────────────────────────
\section{Datasets}

\subsection{MAVEN}

MAVEN~\cite{wang2020maven} is a large-scale event detection dataset covering 168 event types from Wikipedia documents. Each instance is a sentence containing one or more event triggers. We use the training split for evaluation, randomly sampling 50 instances.

\textbf{Why MAVEN may understate model capability:} MAVEN is primarily an event detection benchmark and does not include full argument structure. Many sentences contain multiple plausible triggers, making exact-match evaluation difficult. We use MAVEN as a starting baseline and transition to WikiEvents for richer evaluation.

\subsection{WikiEvents}

WikiEvents~\cite{li2021document} is the standard benchmark for document-level event extraction from the NAACL 2021 paper \textit{Document-Level Event Argument Extraction by Conditional Generation}. It contains 246 annotated documents across a three-level ontology with 49 unique event types. We evaluate on the development split (20 documents, 345 event mentions).

Key statistics for the development split are shown in Table~\ref{tab:wikievents_stats}.

\begin{table}[h]
\centering
\caption{WikiEvents development split statistics.}
\label{tab:wikievents_stats}
\begin{tabular}{lr}
\toprule
\textbf{Statistic} & \textbf{Value} \\
\midrule
Documents & 20 \\
Event mentions & 345 \\
Events with arguments & 291 (84.5\%) \\
Total arguments & 484 \\
Unique full event types & 49 \\
Top-level categories & 13 \\
\bottomrule
\end{tabular}
\end{table}

WikiEvents event types follow a three-level hierarchy (e.g.\ \texttt{Conflict.Attack.DetonateExplode}), making type classification substantially harder than MAVEN's flat label space.

% ─────────────────────────────────────────────────────────────────────────────
\section{Methodology}

\subsection{Models}

We evaluate two open-weight instruction-tuned models:

\textbf{Qwen2.5-7B-Instruct}~\cite{qwen2025}: 7 billion parameters, trained by Alibaba Cloud with strong multilingual instruction following.

\textbf{Llama-3.1-8B-Instruct}~\cite{llama3}: 8 billion parameters, trained by Meta AI with a different pretraining corpus and instruction tuning pipeline.

Both models run on a single NVIDIA A100-SXM4-80GB GPU (Sheffield Stanage HPC) in \texttt{float16} precision with greedy decoding. No fine-tuning is performed on either model.

\subsection{Task Formulation}

For each event mention, we extract the containing sentence and prompt the model to output a JSON object with two fields:
\begin{itemize}
    \item \texttt{trigger}: the word or phrase most clearly expressing the event
    \item \texttt{event\_type}: the event type label (flat for MAVEN; three-level hierarchy for WikiEvents)
\end{itemize}

\subsection{Prompting Strategies}

We compare four prompting strategies on WikiEvents and two on MAVEN:

\textbf{Unconstrained zero-shot:} The prompt explains the output format and shows example type hierarchies but does not enumerate valid labels. The model must generate event types freely.

\textbf{Constrained-label zero-shot:} The prompt explicitly lists all valid event type labels (49 for WikiEvents, 168 for MAVEN). The model is instructed to select only from this list.

\textbf{Few-shot constrained:} Extends the constrained prompt with three worked examples drawn from the training split, covering Conflict, Life, and Contact event categories. Examples demonstrate the expected trigger and type format.

\textbf{Rule-based majority-class baseline:} Always predicts the most frequent event type from the training split (\texttt{Conflict.Attack.DetonateExplode}) and the most frequent trigger word (\texttt{killed}). This establishes a lower-bound floor against which LLM results are compared.

\subsection{Evaluation}

We use a multi-level evaluation framework covering exact match, fuzzy match, and F1 scoring.

\textbf{Exact match} is the strictest metric: the predicted string must exactly equal the gold label (case-insensitive). We report exact match for both trigger and type slots.

\textbf{Fuzzy match (Levenshtein similarity)} measures approximate string closeness. The normalised Levenshtein similarity between strings $s_1$ and $s_2$ is:
\[
\text{sim}(s_1, s_2) = 1 - \frac{\text{lev}(s_1, s_2)}{\max(|s_1|, |s_2|)}
\]
where $\text{lev}$ is the edit distance (number of character insertions, deletions, or substitutions). A fuzzy match is declared when similarity $\geq 0.8$. This captures near-misses such as \textit{killed} vs \textit{kills} (similarity 0.83) that exact match penalises unfairly.

\textbf{Hierarchical partial credit} applies to WikiEvents' three-level type hierarchy. Rather than binary correctness, we award:
\begin{itemize}
    \item 1.0 — exact match (\texttt{Life.Die.Unspecified} = \texttt{Life.Die.Unspecified})
    \item 0.67 — correct Category and SubType (\texttt{Life.Die.*} predicted as \texttt{Life.Die.Unspecified})
    \item 0.33 — correct Category only (\texttt{Life.*} predicted as \texttt{Life.Die.Unspecified})
    \item 0.0 — wrong Category entirely
\end{itemize}

\textbf{Precision, Recall, and F1} are computed per slot. For trigger and type independently:
\[
\text{Precision} = \frac{TP}{TP+FP}, \quad \text{Recall} = \frac{TP}{TP+FN}, \quad F_1 = \frac{2 \cdot P \cdot R}{P + R}
\]

\textbf{Micro F1} aggregates TP, FP, FN counts across all instances before computing F1 — giving more weight to frequent event types.

\textbf{Macro F1} computes F1 per event type class, then averages across all classes — treating rare and frequent types equally. Macro F1 is the more informative metric when the type distribution is skewed.

\textbf{Slot-level evaluation template:} Every prediction is recorded as a structured row with fields for gold and predicted values, Levenshtein similarity, fuzzy match flag, partial credit score, and exact match flag for each slot. Full templates are saved as TSV files in the results directory.

We deliberately avoid BERTScore. Event extraction requires ontology-correct labels — a prediction of \texttt{Life.Injure} when the gold is \texttt{Life.Die} is wrong regardless of semantic proximity, and semantic similarity would unfairly reward such errors.

% ─────────────────────────────────────────────────────────────────────────────
\section{Results}

\subsection{Main Results}

Table~\ref{tab:main_results} presents the full results across all six experimental conditions using exact match, fuzzy match, and F1 metrics.

\begin{table}[h]
\centering
\caption{Main results across all experiments. Trig = trigger slot, Type = event type slot.}
\label{tab:main_results}
\begin{tabular}{llccccccc}
\toprule
\textbf{Model} & \textbf{Condition} & \textbf{N} & \textbf{Trig Exact} & \textbf{Trig Fuzzy} & \textbf{Type Exact} & \textbf{Type PC} & \textbf{Type Macro F1} & \textbf{Both F1} \\
\midrule
\multicolumn{9}{l}{\textit{MAVEN}} \\
Qwen-7B & Unconstrained  & 50  & 0.160 & 0.160 & 0.000 & 0.000 & 0.000 & 0.000 \\
Qwen-7B & Constrained    & 50  & 0.220 & 0.220 & 0.280 & 0.280 & 0.170 & 0.100 \\
\midrule
\multicolumn{9}{l}{\textit{WikiEvents}} \\
Rule-based & Majority class & 345 & — & — & — & — & — & — \\
Qwen-7B & Unconstrained  & 345 & 0.272 & 0.281 & 0.128 & 0.200 & 0.043 & 0.067 \\
Qwen-7B & Constrained    & 345 & 0.264 & 0.272 & 0.186 & 0.279 & 0.079 & 0.113 \\
Qwen-7B & Few-shot       & 345 & 0.275 & 0.278 & 0.252 & 0.333 & 0.164 & 0.157 \\
Llama-3.1-8B & Constrained & 345 & 0.330 & 0.348 & 0.293 & 0.367 & 0.143 & 0.197 \\
\bottomrule
\end{tabular}
\end{table}

PC = hierarchical partial credit (0.33/0.67/1.0). Macro F1 averages F1 per event type class. Rule-based baseline results are shown in Table~\ref{tab:rule_baseline}.

\begin{table}[h]
\centering
\caption{Rule-based majority-class baseline on WikiEvents dev ($n=345$).}
\label{tab:rule_baseline}
\begin{tabular}{lccc}
\toprule
\textbf{Strategy} & \textbf{Trigger Exact} & \textbf{Type Exact} & \textbf{Both} \\
\midrule
Always predict \texttt{killed} / \texttt{Conflict.Attack.DetonateExplode} & 0.122 & 0.191 & 0.000 \\
\bottomrule
\end{tabular}
\end{table}

The rule-based baseline achieves non-trivial type accuracy (0.191) purely from class frequency, but zero joint accuracy. Every LLM condition substantially exceeds this on the combined trigger+type metric, confirming that LLMs add genuine extraction capability beyond majority-class prediction.

\subsection{Effect of Prompting Strategy}

Three prompting effects are visible in the WikiEvents results:

\textbf{Constraining labels improves type accuracy} (unconstrained 0.128 → constrained 0.186). The jump reflects reduced ontology vocabulary mismatch — without constraints, Qwen produces semantically plausible but benchmark-incompatible type strings. Label adherence improved from free-form generation to 80.6\% within the valid label set under constrained prompting.

\textbf{Few-shot examples further improve type accuracy} (constrained 0.186 → few-shot 0.252, +35\% relative). More strikingly, label set adherence jumps from 80.6\% to 97.4\% with examples — showing that worked examples teach the model to follow the ontology format more reliably than the label list alone.

\textbf{Trigger accuracy is largely unaffected by prompting strategy} (range 0.264--0.275 across Qwen WikiEvents conditions). This confirms that trigger selection and type classification are independent problems: prompt constraints help ontology alignment but not event selection under ambiguity.

\subsection{Model Comparison: Qwen-7B vs.\ Llama-3.1-8B}

Llama-3.1-8B outperforms Qwen-7B on every WikiEvents metric under matched conditions (both constrained zero-shot):

\begin{itemize}
    \item Trigger exact: 0.264 → 0.330 (+25\% relative)
    \item Type exact: 0.186 → 0.293 (+57\% relative)
    \item Type partial credit: 0.279 → 0.367 (+32\% relative)
    \item Macro F1: 0.079 → 0.143 (+81\% relative)
    \item Label adherence: 80.6\% → 98.8\%
\end{itemize}

The label adherence gap is particularly notable — Llama follows the constrained label list far more reliably than Qwen under the same prompt. This suggests that Llama's instruction tuning produces stronger constraint-following behaviour. Despite identical parameter count ($\approx$8B), the two models differ substantially in capability for this task.

\subsection{Micro F1 vs.\ Macro F1}

The gap between Micro and Macro F1 reveals class imbalance effects. WikiEvents unconstrained Qwen achieves Micro F1 = 0.128 but Macro F1 = 0.043. The model performs well on the most frequent types (\texttt{Conflict.Attack}: 0.439--0.472) but scores zero on rare types, pulling the macro average down sharply. Macro F1 is the more honest measure of generalisation across the ontology.

\subsection{Fuzzy Match vs.\ Exact Match}

Fuzzy trigger match provides a small but consistent gain over exact match (e.g.\ Llama: 0.330 exact vs.\ 0.348 fuzzy), indicating some near-miss predictions where the model predicts a morphological variant of the correct trigger (e.g.\ \textit{kills} vs.\ \textit{killed}). For type matching, fuzzy and exact scores are nearly identical across all conditions — type errors are categorical wrong choices, not spelling near-misses, confirming that the primary challenge is ontology alignment rather than string formatting.

% ─────────────────────────────────────────────────────────────────────────────
\section{Error Analysis}

\subsection{MAVEN Error Analysis}

Table~\ref{tab:maven_errors} shows the error breakdown for both MAVEN conditions.

\begin{table}[h]
\centering
\caption{Error category breakdown — MAVEN ($n=50$ per condition).}
\label{tab:maven_errors}
\begin{tabular}{lrrrr}
\toprule
\textbf{Error Category} & \multicolumn{2}{c}{\textbf{Unconstrained}} & \multicolumn{2}{c}{\textbf{Constrained}} \\
\cmidrule(lr){2-3}\cmidrule(lr){4-5}
 & Count & \% & Count & \% \\
\midrule
Correct (trigger + type)     &  0 &  0.0\% &  5 & 10.0\% \\
Right trigger, wrong type    &  8 & 16.0\% &  6 & 12.0\% \\
Wrong trigger, right type    &  0 &  0.0\% &  9 & 18.0\% \\
Wrong trigger, wrong type    & 42 & 84.0\% & 30 & 60.0\% \\
\bottomrule
\end{tabular}
\end{table}

The most striking finding in the unconstrained condition is that type accuracy is exactly 0.000 — not because the model fails to understand event semantics, but because it \textit{never uses MAVEN's ontology labels}. Instead it generates its own vocabulary: \texttt{natural\_disaster}, \texttt{Natural Disaster}, \texttt{tectonic}, and in two cases the Chinese equivalent \texttt{自然灾害}. Without knowledge of the exact label strings, the model produces semantically reasonable but benchmark-incompatible outputs.

Constraining the label space transforms the error distribution substantially. The ``wrong trigger, right type'' category jumps from 0.0\% to 18.0\%, indicating that label constraints successfully guide type classification even in cases where trigger selection fails. The dominant error category (wrong trigger, wrong type) drops from 84.0\% to 60.0\%, confirming that much of MAVEN's difficulty in the unconstrained condition is ontology vocabulary mismatch rather than genuine semantic failure.

The top type confusions in the constrained condition reflect the domain of the evaluated sample (earthquake/tsunami events): \texttt{Causation} $\rightarrow$ \texttt{Catastrophe} (3 cases) and \texttt{Destroying} $\rightarrow$ \texttt{Catastrophe} (2 cases) are both plausible labellings of the same physical events, highlighting the difficulty of MAVEN's fine-grained ontology distinctions.

\subsection{WikiEvents Error Categories}

Table~\ref{tab:errors} categorises all predictions from the WikiEvents constrained condition. Errors are classified by whether the trigger word and each level of the event type hierarchy are correct.

\begin{table}[h]
\centering
\caption{Error category breakdown — WikiEvents constrained zero-shot ($n=345$).}
\label{tab:errors}
\begin{tabular}{lrr}
\toprule
\textbf{Error Category} & \textbf{Count} & \textbf{\%} \\
\midrule
Correct (trigger + type)                        & 39  & 11.3\% \\
Right trigger, right Cat.Sub, wrong Specificity & 16  &  4.6\% \\
Right trigger, right Category, wrong SubType    & 13  &  3.8\% \\
Right trigger, wrong Category entirely          & 23  &  6.7\% \\
Wrong trigger, right type                       & 25  &  7.2\% \\
Wrong trigger, right Cat.Sub                    & 15  &  4.3\% \\
Wrong trigger, right Category, wrong SubType    & 22  &  6.4\% \\
Wrong trigger, wrong Category entirely          & 192 & 55.7\% \\
\bottomrule
\end{tabular}
\end{table}

The dominant error category is \textit{wrong trigger, wrong category entirely} (55.7\%), indicating that the model's primary failure mode is not subtle ontology misalignment but rather selecting a fundamentally different event within the same sentence.

\subsection{Effect of Constraining on Error Distribution}

Compared to unconstrained prompting, constraining labels halves the ``right trigger, wrong category'' error rate (15.4\% → 6.7\%). This confirms that label constraints provide meaningful ontology alignment. However, the most frequent error category (wrong trigger + wrong category) shifts only marginally (57.4\% → 55.7\%), indicating that this core failure is driven by trigger selection rather than type knowledge.

\subsection{Type Confusions}

The most frequent type confusions reveal a systematic pattern: the model conflates \textit{consequence} events with \textit{cause} events. Table~\ref{tab:confusions} shows the top confusions.

\begin{table}[h]
\centering
\caption{Top type confusions (gold $\rightarrow$ predicted), WikiEvents constrained.}
\label{tab:confusions}
\begin{tabular}{llr}
\toprule
\textbf{Gold Type} & \textbf{Predicted Type} & \textbf{Count} \\
\midrule
Life.Die.Unspecified       & Conflict.Attack.DetonateExplode & 14 \\
Life.Die.Unspecified       & Conflict.Attack.Unspecified     & 12 \\
Life.Injure.Unspecified    & Life.Injure                     &  8 \\
Life.Die.Unspecified       & Life.Injure                     &  7 \\
Conflict.Attack.DetonateExplode & Conflict.Attack.Unspecified &  6 \\
\bottomrule
\end{tabular}
\end{table}

The most frequent confusion (Life.Die → Conflict.Attack.DetonateExplode, 14 cases) is explained by multi-event sentences: when a bombing kills someone, both events are linguistically present. The model attaches to the more physically salient explosion event rather than the death consequence. Eight cases of Life.Injure.Unspecified → Life.Injure indicate a \textit{format failure} rather than a semantic failure: the model predicts the correct category and subtype but drops the required third level of the hierarchy.

\subsection{Qualitative Examples}

\textbf{Multi-event confusion (dominant failure):}
\begin{quote}
\small
\textit{Sentence:} ``Roadside IED kills Russian major general in Syria''\\
\textit{Gold:} trigger=\texttt{kills}, type=\texttt{Life.Die.Unspecified}\\
\textit{Predicted:} trigger=\texttt{kills}, type=\texttt{Conflict.Attack.DetonateExplode}
\end{quote}
The trigger is correct, but the model classifies the overall event as an attack rather than a death — plausible given the sentence content but ontologically wrong.

\textbf{Trigger selection under ambiguity:}
\begin{quote}
\small
\textit{Sentence:} ``Russia's Defense Ministry confirmed that a major general was killed in Syria by an IED...''\\
\textit{Gold:} trigger=\texttt{killed}, type=\texttt{Life.Die.Unspecified}\\
\textit{Predicted:} trigger=\texttt{confirmed}, type=\texttt{Contact.Contact.Broadcast}
\end{quote}
The model selects \texttt{confirmed} as the trigger, anchoring on the reporting verb rather than the actual event being described — a common failure mode in news-domain text.

\textbf{Correct prediction:}
\begin{quote}
\small
\textit{Sentence:} ``A local commander of Syria's National Defense Forces was also reportedly killed.''\\
\textit{Gold:} trigger=\texttt{killed}, type=\texttt{Life.Die.Unspecified}\\
\textit{Predicted:} trigger=\texttt{killed}, type=\texttt{Life.Die.Unspecified} ✓
\end{quote}
Single-event sentences with unambiguous triggers are reliably extracted.

% ─────────────────────────────────────────────────────────────────────────────
\section{Discussion}

\subsection{What Zero-Shot LLMs Can Do}

Our results demonstrate that Qwen2.5-7B-Instruct can perform meaningful zero-shot event extraction. JSON output validity is perfect across all experiments (345/345), showing reliable instruction-following. Trigger accuracy reaches 0.272 on WikiEvents without any examples, comparable to weak supervised baselines from earlier work. The model shows genuine understanding of event semantics, as evidenced by the semantic proximity of its trigger errors (e.g.\ \textit{injured} → \textit{killed}).

\subsection{Where Zero-Shot LLMs Struggle}

Three failure modes account for most errors:

\textbf{Multi-event sentence ambiguity.} News text routinely describes multiple co-occurring events. The model consistently selects the most salient event in the sentence rather than the gold-annotated one. This is an inherent limitation of sentence-level prompting for document-level annotations.

\textbf{Ontology misalignment.} Even with constrained prompting, 19.4\% of WikiEvents predictions used types outside the provided label list. The three-level hierarchy is harder to constrain than a flat label space.

\textbf{Hierarchical format adherence.} Eight cases of Life.Injure.Unspecified → Life.Injure suggest the model drops the third hierarchy level, indicating that longer, structured label strings require more explicit format enforcement in the prompt.

\subsection{Implications}

These results establish a clear zero-shot baseline against which fine-tuned or few-shot approaches can be compared. The gap between unconstrained and constrained accuracy (particularly on type classification) suggests that prompt engineering alone can recover substantial performance. However, the dominance of the multi-event confusion error class (55.7\%) suggests that sentence-level prompting is fundamentally limited for document-annotated benchmarks like WikiEvents.

% ─────────────────────────────────────────────────────────────────────────────
\section{Conclusion}

We presented a zero-shot LLM baseline for event extraction using Qwen2.5-7B-Instruct on MAVEN and WikiEvents. Key findings:

\begin{itemize}
    \item Constrained-label prompting consistently improves type accuracy (MAVEN: 0.000 → 0.280; WikiEvents: 0.128 → 0.186)
    \item Trigger accuracy is stronger on WikiEvents (0.272) than MAVEN (0.160), likely due to more salient event types
    \item The dominant failure mode is multi-event sentence ambiguity, accounting for 55.7\% of all errors
    \item Hierarchical ontologies are harder to constrain than flat label spaces, with 19.4\% of predictions falling outside the provided label set
\end{itemize}

Future work could explore few-shot prompting, document-level context, or fine-tuned models to address the multi-event ambiguity problem identified here.

% ─────────────────────────────────────────────────────────────────────────────
\bibliographystyle{plain}
\begin{thebibliography}{99}

\bibitem{ahn2006stages}
D.\ Ahn.
\newblock The stages of event extraction.
\newblock In \textit{Proceedings of the Workshop on Annotating and Reasoning about Time and Events}, 2006.

\bibitem{li2021document}
S.\ Li, H.\ Ji, and J.\ Han.
\newblock Document-level event argument extraction by conditional generation.
\newblock In \textit{Proceedings of NAACL}, 2021.

\bibitem{li2013joint}
Q.\ Li and H.\ Ji.
\newblock Incremental joint extraction of entity mentions and relations.
\newblock In \textit{Proceedings of ACL}, 2013.

\bibitem{lu2021text2event}
Y.\ Lu et al.
\newblock Text2Event: Controllable sequence-to-structure generation for end-to-end event extraction.
\newblock In \textit{Proceedings of ACL}, 2021.

\bibitem{ma2023large}
Y.\ Ma et al.
\newblock Large language model is not a good few-shot information extractor, but a good reranker for hard samples!
\newblock In \textit{Findings of EMNLP}, 2023.

\bibitem{nguyen2015event}
T.\ H.\ Nguyen and R.\ Grishman.
\newblock Event detection and domain adaptation with convolutional neural networks.
\newblock In \textit{Proceedings of ACL}, 2015.

\bibitem{llama3}
A.\ Dubey et al.
\newblock The Llama 3 herd of models.
\newblock \textit{arXiv preprint arXiv:2407.21783}, 2024.

\bibitem{qwen2025}
Qwen Team.
\newblock Qwen2.5: A party of foundation models.
\newblock \textit{arXiv preprint arXiv:2412.15115}, 2025.

\bibitem{wang2020maven}
X.\ Wang et al.
\newblock MAVEN: A massive general domain event detection dataset.
\newblock In \textit{Proceedings of EMNLP}, 2020.

\bibitem{wei2023zero}
X.\ Wei et al.
\newblock Zero-shot information extraction via chatting with ChatGPT.
\newblock \textit{arXiv preprint arXiv:2302.10205}, 2023.

\bibitem{yang2019exploring}
S.\ Yang and K.\ Mitchel.
\newblock Exploring pre-trained language models for event extraction and generation.
\newblock In \textit{Proceedings of ACL}, 2019.

\end{thebibliography}

\end{document}
